\chapter*{Remerciements}
\addcontentsline{toc}{chapter}{Remerciements}

Je tiens tout d'abord à remercier très sincèrement \textbf{Yassmina Sadi}, mon encadrante à l'Université d'Évry Paris-Saclay, pour sa confiance, sa disponibilité et la qualité de son accompagnement tout au long de ce stage. Ses conseils, ses relectures attentives et ses exigences m'ont aidé à structurer ma démarche, à prendre du recul sur mes résultats et à aller au bout de chaque question.

J'adresse également mes sincères remerciements à \textbf{Lamjed Ben Said}, qui a encadré ce travail à l'Institut Supérieur de Gestion de Tunis, pour son suivi, ses conseils et la confiance qu'il m'a accordée tout au long de ce mémoire.

Je remercie chaleureusement le \textbf{Pr Nicholas Heming}, professeur de médecine intensive--réanimation à l'Hôpital Raymond-Poincaré (AP-HP), pour le temps considérable qu'il a consacré à la validation médicale du guideline et du Gold Standard. Ses explications cliniques et sa rigueur ont donné à ce travail une base solide, sans laquelle aucune évaluation fiable n'aurait été possible.

J'adresse également mes remerciements à l'\textbf{IHU SEPSIS} et à l'ensemble des membres du Work Package~6, qui m'ont accueilli et permis de travailler sur les dossiers de la cohorte APROCCHSS. Je remercie en particulier \textbf{Eya Ben Dhiaf}, dont les travaux sur le graphe de connaissances ont servi de point de départ à ce mémoire, pour nos échanges, son aide et notre collaboration lors des présentations de ce projet.

Je remercie les membres du jury, qui ont accepté d'évaluer ce travail et d'y consacrer leur temps.

Mes remerciements vont aussi à l'ensemble des enseignants du \textbf{master IDEAS (ID2M)} de l'\textbf{Institut Supérieur de Gestion de Tunis}, pour la formation qu'ils m'ont apportée et qui m'a permis de mener ce projet. Je remercie également l'\textbf{Université d'Évry Paris-Saclay}, qui m'a accueilli pour ce stage.

Enfin, je remercie du fond du cœur ma famille et mes proches, pour leur soutien constant, leur patience et leurs encouragements tout au long de mes études. Ce travail leur doit beaucoup.
